\chapter{Totalität bei Mengenoperationen}
\label{b04guide:ch:sets}

\section{Einschränken und Erweitern}

\begin{definition}[Einschränkung auf Ausgangs- und Zielmengen]
Für $R\subseteq A\times B$, $C\subseteq A$ und $D\subseteq B$ sei
\[
 R_{C,D}=R\cap(C\times D).
\]

\end{definition}

\begin{theorem}[Einschränkung der Ausgangsmenge]
\label{b04guide:thm:restrict}
Für $\TR(R;A,B)$ und $C\subseteq A$ gilt
\[
 \TR(R_{C,B};C,B).
\]

\end{theorem}

\begin{theorem}[Genaues Kriterium für eine Zielbeschränkung]
\label{b04guide:thm:target}
Für $\TR(R;A,B)$, $C\subseteq A$ und $D\subseteq B$ gilt
\[
 \TR(R_{C,D};C,D)
 \quad\Longleftrightarrow\quad
 \forall a\in C\,(R_a\cap D\ne\varnothing).
\]
Die eingeschränkte Zielmenge muß also jede der betreffenden
Fasern treffen. Allein $D\ne\varnothing$ genügt im allgemeinen nicht.
\end{theorem}

\begin{theorem}[Hinzufügen von Paaren]\label{b04guide:thm:upward}
Für $\TR(R;A,B)$ und $R\subseteq S\subseteq A\times B$ gilt
\[
 \TR(S;A,B).
\]

\end{theorem}

\begin{theorem}[Erweiterung der Ausgangsmenge]
\label{b04guide:thm:extend}
Seien $\TR(R;A,B)$, $A\subseteq A'$ und $B\subseteq B'$.
Dann existiert eine totale Erweiterung genau unter folgender Bedingung:
\[
 \begin{gathered}
 \exists S\in\TRset(A',B')\,(R\subseteq S)\\
 \Longleftrightarrow\quad A'=A\ \lor\ B'\ne\varnothing.
 \end{gathered}
\]
Ist diese Bedingung erfüllt, so ist insbesondere
\[
 R\cup\bigl((A'\setminus A)\times B'\bigr)
\]
eine solche Erweiterung. Ihre Einschränkung auf $A\times B'$ ist
genau $R$.
\end{theorem}

\section{Vereinigen und Zusammenfügen}

\begin{theorem}[Nichtleere Vereinigungen totaler Relationen]
\label{b04guide:thm:union}
Für $\varnothing\ne\mathcal R\subseteq\TRset(A,B)$ gilt
\[
 \TR\bigl(\bigcup\mathcal R;A,B\bigr).
\]
Für $\mathcal R=\varnothing$ ist die Vereinigung genau dann
total von $A$ nach $B$, wenn $A=\varnothing$.
Insbesondere ist $R\cup S$ total, sobald $R$ total ist und
$S\subseteq A\times B$ gilt.
\end{theorem}

\begin{theorem}[Zusammenfügen auf einer Überdeckung]
\label{b04guide:thm:gluing}
Seien $R,S$ totale Relationen mit
\[
 \TR(R;A,B),\qquad \TR(S;C,B).
\]
Dann gilt
\[
 \TR(R\cup S;A\cup C,B).
\]
Auf $A\cap C$ dürfen $R$ und $S$ unterschiedliche Paare enthalten.
Es ist keine Übereinstimmungsbedingung erforderlich.
\end{theorem}

\begin{theorem}[Mengenweise Zusammenfügen]
\label{b04guide:thm:familyglue}
Seien $\mathcal C\subseteq\mathcal P(A)$ mit
$\bigcup\mathcal C=A$ und
$\mathcal R\subseteq\mathcal P(A\times B)$.
Wenn
\[
 \forall C\in\mathcal C\,\exists R\in\mathcal R\,
                  \TR(R;C,B),
\]
dann gilt
\[
 \TR\bigl(\bigcup\mathcal R;A,B\bigr).
\]
Hier werden sämtliche vorliegenden Paare vereinigt.
Eine Auswahl eines einzelnen $R$ für jedes $C$ ist kein
Bestandteil der Aussage.
\end{theorem}

\begin{theorem}[Ersetzen auf einer Teilmenge]
\label{b04guide:thm:override}
Für $\TR(R;A,B)$, $C\subseteq A$ und $\TR(S;C,B)$ gilt
\[
 \TR\bigl((R\setminus(C\times B))\cup S;A,B\bigr).
\]
Die neue Relation stimmt außerhalb von $C$ mit $R$ und
auf $C$ mit $S$ überein.
\end{theorem}

\section{Durchschnitt und Entfernen von Paaren}

\begin{theorem}[Durchschnitt zweier totaler Relationen]
\label{b04guide:thm:intersection}
Für $\TR(R;A,B)$ und $\TR(S;A,B)$ gilt
\[
 \TR(R\cap S;A,B)
 \quad\Longleftrightarrow\quad
 \forall a\in A\,(R_a\cap S_a\ne\varnothing).
\]
Der Durchschnitt totaler Relationen ist im allgemeinen nicht
total: Sind $A\ne\varnothing$ und $b,c\in B$ mit $b\ne c$,
so sind $A\times\{b\}$ und $A\times\{c\}$ total,
ihr Durchschnitt jedoch nicht.
\end{theorem}

\begin{definition}[Relativer Durchschnitt einer Relationenmenge]
Für $\mathcal R\subseteq\mathcal P(A\times B)$ sei
\[
 I_{A,B}(\mathcal R)
 =\{(a,b)\in A\times B\mid
                         \forall R\in\mathcal R\,((a,b)\in R)\}.
\]
Insbesondere gilt $I_{A,B}(\varnothing)=A\times B$.
Für $\mathcal R\ne\varnothing$ ist dies der übliche
Mengendurchschnitt $\bigcap\mathcal R$.
\end{definition}

\begin{theorem}[Durchschnitt einer Familie totaler Relationen]
\label{b04guide:thm:familyintersection}
Für $\mathcal R\subseteq\TRset(A,B)$ gilt
\[
 \begin{gathered}
 \TR\bigl(I_{A,B}(\mathcal R);A,B\bigr)\\
 \Longleftrightarrow\
 \forall a\in A\,\exists b\in B\,
                      \forall R\in\mathcal R\,(b\in R_a).
 \end{gathered}
\]
Dies gilt auch für die leere Familie.
Hat $\mathcal R$ eine gemeinsame totale Teilrelation, also ein
$T\in\TRset(A,B)$ mit $T\subseteq R$ für alle
$R\in\mathcal R$, so ist der relative Durchschnitt total.
\end{theorem}

\begin{theorem}[Entfernen beliebiger Paare]
\label{b04guide:thm:difference}
Für $\TR(R;A,B)$ und $E\subseteq A\times B$ gilt
\[
 \TR(R\setminus E;A,B)
 \quad\Longleftrightarrow\quad
 \forall a\in A\,(R_a\setminus E_a\ne\varnothing).
\]
Für $(a,b)\in R$ gilt insbesondere
\[
 \TR\bigl(R\setminus\{(a,b)\};A,B\bigr)
 \quad\Longleftrightarrow\quad R_a\setminus\{b\}\ne\varnothing.
\]
Ein Paar darf genau dann entfernt werden, wenn über seiner
ersten Komponente noch ein weiteres Paar erhalten bleibt.
\end{theorem}

\begin{theorem}[Komplement einer totalen Relation]
\label{b04guide:thm:complement}
Für $\TR(R;A,B)$ gilt
\[
 \TR\bigl((A\times B)\setminus R;A,B\bigr)
 \quad\Longleftrightarrow\quad
 \forall a\in A\,(R_a\ne B).
\]
Damit sind $R$ und sein relatives Komplement genau dann beide
total, wenn für jedes $a\in A$ gilt
\[
 \varnothing\ne R_a\subsetneq B.
\]
\end{theorem}
